\documentclass[ngerman]{beamer}
% Option handout um Version ohne \pause etc. zu erzeugen

\usepackage[utf8]{inputenc}
\usepackage[T1]{fontenc}
\usepackage{babel}

\usepackage{booktabs}
\usepackage{graphicx}
\usepackage{csquotes}
\usepackage{xcolor}

%\usetheme{PaloAlto}

\author{Uwe Ziegenhagen}
\title{Meine erste Präsentation}
\subtitle{Meine Abenteuer in Entenhausen}
\date{Köln, den \today}

\begin{document}

\begin{frame}

\maketitle

\end{frame}

\begin{frame}
\frametitle{Unser heutiges Programm}
\framesubtitle{Tag 1}

\tableofcontents

\end{frame}





\section{Grundlagen}

\frame{
\frametitle{Hallo Fernuni Hagen!}

\begin{itemize}
	\item Hallo
	\item ich 
	\item bin 
	\item eine 
	\item Bullet
	\item Liste
\end{itemize}

}


\frame{
\frametitle{Hallo Fernuni Hagen!}

\begin{enumerate}
	\item Hallo
	\item ich 
	\item bin 
	\item eine 
	\begin{enumerate}
	\item Hallo
	\item ich 
	\item bin 
	\item eine 
	\item Bullet
	\item Liste
\end{enumerate}
	\item Bullet
	\item Liste
\end{enumerate}

}


\begin{frame}
\frametitle{Bilder gehen auch}

\begin{figure}[htb]
	\centering
	\includegraphics[width=\textwidth]{miau.jpg}
	\caption{Das ist eine Katze}\label{fig:mieze2}
\end{figure}

\end{frame}


\begin{frame}
\frametitle{Zweispaltig}

\begin{columns}

\begin{column}{0.33\textwidth}
	\begin{itemize}
		\item Ich bin 
		\item eine 
		\item Bullet-Liste
	\end{itemize}
\end{column}

\begin{column}{0.66\textwidth}
	\begin{figure}[htb]
		\centering
		\includegraphics[width=0.9\textwidth]{miau.jpg}
		\caption{Das ist eine Katze}\label{fig:mieze2}
	\end{figure}
\end{column}

\end{columns}
\end{frame}

\begin{frame}
\frametitle{Mathematik}

\[a^2 + b^2 = c^2 \]

\begin{equation}
a^2 + b^2 = c^2
\end{equation}


\end{frame}

\begin{frame}

\begin{itemize}
	\item Hallo \pause
	\item ich \pause
	\item bin  \pause
	\item eine \pause 
	\item Bullet
	\item Liste
\end{itemize}

\end{frame}

\section{Aufdeckungen}

\begin{frame}

\begin{itemize}
	\item<1-> Hallo 
	\item<2> ich 
	\item<-2> bin  
	\item<3-> eine 
	\item<1> Bullet
	\item<2> Liste
\end{itemize}

\end{frame}


\begin{frame}

\begin{itemize}
	\item<1-> Hallo 
	\item<2-> ich 
	\item<3-> bin  
	\item<4-> eine 
	\item<5-> Bullet
	\item<6-> Liste
\end{itemize}

\end{frame}


\end{document}